% GENERATED FILE. Edit canonical lessons, then run npm run generate:books.
% canonical-source-hash: fnv1a64:0a18f63cbe1e724a
% canonical-lessons: RU-C79-bus, RU-C79-train, RU-C79-plane, RU-C79-car, RU-C79-bridge

\chapter{Bus, Train, Plane, Car, Bridge}
\label{ch:ru-bus-train-plane-car-bridge}
\hlchaptermodality{79}

\begin{chapteropening}
\textbf{By the end of this chapter:} \emph{I can name a bus, a train, a plane, a car and a bridge.}

Complete the last lesson of chapter 79: I can name a bus, a train, a plane, a car and a bridge.
\end{chapteropening}

\section[\texorpdfstring{\ru{автобус} (avtóbus)}{автобус (avtóbus)}]{\ru{автобус} (avtóbus) — a bus}

\label{lesson:RU-C79-bus}



\begin{quote}
Before the new one: say the Russian for a theatre, then the Russian for a restaurant.
\end{quote}

\subsection*{You'll want to know: \ru{автобус}}
\textbf{\ru{автобус}} (\emph{avtóbus}) — \textquotedblleft{}a bus\textquotedblright{}.

\begin{cousinweb}
\textbf{\ru{автобус}} is an international word of the early 1900s: \emph{\ru{авто}-}, from Greek \emph{autós}, \textquotedblleft{}self\textquotedblright{}, as in \emph{automobile}, with \emph{\ru{бус}}, cut from Latin \emph{omnibus}, \textquotedblleft{}for everyone\textquotedblright{}.
\end{cousinweb}

\begin{grammarlens}[title={What you've built}]
One more word for this chapter.
\end{grammarlens}

\subsection*{Your turn}
\begin{itemize}
\raggedright
  \item \emph{Say it:} \emph{avtóbus}
  \item \emph{Say it:} \emph{avtóbus}, once more
  \item \emph{Say it:} say \emph{avtóbus}
  \item \emph{Recall:} say the Russian for a theatre, then the Russian for a restaurant, then say \emph{avtóbus} again
\end{itemize}

\begin{tcolorbox}[breakable,colback=teal!4,colframe=teal!35!black,title={Before you move on}]
What does \ru{автобус} mean? (\textquotedblleft{}A bus\textquotedblright{}.) Say it once more.
\end{tcolorbox}
\section[\texorpdfstring{\ru{поезд} (póyezd)}{поезд (póyezd)}]{\ru{поезд} (póyezd) — a train}

\label{lesson:RU-C79-train}



\begin{quote}
Before the new one: say the Russian for a restaurant, then the Russian for a bus.
\end{quote}

\subsection*{You'll want to know: \ru{поезд}}
\textbf{\ru{поезд}} (\emph{póyezd}) — \textquotedblleft{}a train\textquotedblright{}.

\begin{cousinweb}
\textbf{\ru{поезд}} is \emph{\ru{по}-} with \emph{\ru{езд}}, from \emph{\ru{ездить}}, \textquotedblleft{}to ride, to travel\textquotedblright{}, and \emph{\ru{езда}}, \textquotedblleft{}a ride\textquotedblright{} (Vasmer). It first meant \textquotedblleft{}a journey\textquotedblright{}, then \textquotedblleft{}a procession of carriages\textquotedblright{}, then a train.
\end{cousinweb}

\begin{grammarlens}[title={What you've built}]
One more word for this chapter.
\end{grammarlens}

\subsection*{Your turn}
\begin{itemize}
\raggedright
  \item \emph{Say it:} \emph{póyezd}
  \item \emph{Say it:} \emph{póyezd}, once more
  \item \emph{Say it:} say \emph{póyezd}
  \item \emph{Recall:} say the Russian for a restaurant, then the Russian for a bus, then say \emph{póyezd} again
\end{itemize}

\begin{tcolorbox}[breakable,colback=teal!4,colframe=teal!35!black,title={Before you move on}]
What does \ru{поезд} mean? (\textquotedblleft{}A train\textquotedblright{}.) Say it once more.
\end{tcolorbox}
\section[\texorpdfstring{\ru{самолёт} (samolyót)}{самолёт (samolyót)}]{\ru{самолёт} (samolyót) — an aeroplane}

\label{lesson:RU-C79-plane}



\begin{quote}
Before the new one: say the Russian for a bus, then the Russian for a train.
\end{quote}

\subsection*{You'll want to know: \ru{самолёт}}
\textbf{\ru{самолёт}} (\emph{samolyót}) — \textquotedblleft{}an aeroplane\textquotedblright{}. \textquotedblleft{}Self-flier\textquotedblright{}.

\begin{cousinweb}
\textbf{\ru{самолёт}} is \emph{\ru{сам}}, \textquotedblleft{}self\textquotedblright{}, with \emph{-\ru{лёт}}, \textquotedblleft{}flight\textquotedblright{}, from \emph{\ru{летать}}, \textquotedblleft{}to fly\textquotedblright{}: \textquotedblleft{}a self-flier\textquotedblright{}. The fairy-tale flying carpet is the \emph{\ru{ковёр}-\ru{самолёт}}, \textquotedblleft{}carpet-plane\textquotedblright{}.
\end{cousinweb}

\begin{grammarlens}[title={What you've built}]
One more word for this chapter.
\end{grammarlens}

\subsection*{Your turn}
\begin{itemize}
\raggedright
  \item \emph{Say it:} \emph{samolyót}
  \item \emph{Say it:} \emph{samolyót}, once more
  \item \emph{Say it:} say \emph{samolyót}
  \item \emph{Recall:} say the Russian for a bus, then the Russian for a train, then say \emph{samolyót} again
\end{itemize}

\begin{tcolorbox}[breakable,colback=teal!4,colframe=teal!35!black,title={Before you move on}]
What does \ru{самолёт} mean? (\textquotedblleft{}An aeroplane\textquotedblright{}.) Say it once more.
\end{tcolorbox}
\section[\texorpdfstring{\ru{машина} (mashína)}{машина (mashína)}]{\ru{машина} (mashína) — a car}

\label{lesson:RU-C79-car}



\begin{quote}
Before the new one: say the Russian for a train, then the Russian for an aeroplane.
\end{quote}

\subsection*{You'll want to know: \ru{машина}}
\textbf{\ru{машина}} (\emph{mashína}) — \textquotedblleft{}a car\textquotedblright{}.

\begin{cousinweb}
\textbf{\ru{машина}} — from Peter the Great's time; it came through German \emph{Maschine} from French \emph{machine}, Latin \emph{māchina}, \textquotedblleft{}a device\textquotedblright{} (Vasmer). The sense \textquotedblleft{}car\textquotedblright{} is modern.
\end{cousinweb}

\begin{grammarlens}[title={What you've built}]
One more word for this chapter.
\end{grammarlens}

\subsection*{Your turn}
\begin{itemize}
\raggedright
  \item \emph{Say it:} \emph{mashína}
  \item \emph{Say it:} \emph{mashína}, once more
  \item \emph{Say it:} say \emph{mashína}
  \item \emph{Recall:} say the Russian for a train, then the Russian for an aeroplane, then say \emph{mashína} again
\end{itemize}

\begin{tcolorbox}[breakable,colback=teal!4,colframe=teal!35!black,title={Before you move on}]
What does \ru{машина} mean? (\textquotedblleft{}A car\textquotedblright{}.) Say it once more.
\end{tcolorbox}
\section[\texorpdfstring{\ru{мост} (most)}{мост (most)}]{\ru{мост} (most) — a bridge}

\label{lesson:RU-C79-bridge}



\begin{quote}
Before the new one: say the Russian for an aeroplane, then the Russian for a car.
\end{quote}

\subsection*{You'll want to know: \ru{мост}}
\textbf{\ru{мост}} (\emph{most}) — \textquotedblleft{}a bridge\textquotedblright{}.

\begin{cousinweb}
\textbf{\ru{мост}} $\leftarrow$ Proto-Slavic \emph{*mostŭ}, \textquotedblleft{}bridge; a plank floor\textquotedblright{} — inherited. Vasmer finds a Germanic source unreliable but compares Old High German \emph{mast}, \textquotedblleft{}a pole\textquotedblright{}, so the deeper origin is open.
\end{cousinweb}

\begin{grammarlens}[title={What you've built}]
One more word for this chapter.
\end{grammarlens}

\subsection*{Your turn}
\begin{itemize}
\raggedright
  \item \emph{Say it:} \emph{most}
  \item \emph{Say it:} \emph{most}, once more
  \item \emph{Say it:} say \emph{most}
  \item \emph{Recall:} say the Russian for an aeroplane, then the Russian for a car, then say \emph{most} again
\end{itemize}

\begin{tcolorbox}[breakable,colback=teal!4,colframe=teal!35!black,title={Before you move on}]
What does \ru{мост} mean? (\textquotedblleft{}A bridge\textquotedblright{}.) Say it once more.
\end{tcolorbox}
